\documentclass[12pt,a4paper]{article}
\usepackage[utf8]{inputenc}
\usepackage[T1]{fontenc}
\usepackage[norsk]{babel}
\usepackage{geometry}
\geometry{left=2.5cm,right=2.5cm,top=3.5cm,bottom=2.5cm}
\usepackage{graphicx}
\usepackage{fancyhdr}
\usepackage{titlesec}
\usepackage[hidelinks]{hyperref}
\setlength{\headheight}{50pt}
\setlength{\headsep}{20pt}
\pagestyle{fancy}
\fancyhf{}
\fancyhead[L]{\small Skallstuggu Fjellstue -- Personvernerklæring}
% Valgfri logo: last opp logo.png i Overleaf. Dokumentet virker ogsa uten.
\fancyhead[R]{\IfFileExists{logo.png}{\includegraphics[height=1.3cm]{logo.png}}{}}
\fancyfoot[R]{\small Side \thepage}
\renewcommand{\headrulewidth}{0.4pt}
\titleformat{\section}{\large\bfseries}{}{0pt}{}
\titlespacing{\section}{0pt}{12pt}{6pt}
\setlength{\parskip}{6pt}
\setlength{\parindent}{0pt}
\setlength{\emergencystretch}{3em}
\begin{document}
\begin{center}
{\LARGE\textbf{Personvernerklæring}}\\[6pt]
{\large Skallstuggu Fjellstue -- Frolfjellet, Levanger}\\[6pt]
\small Sist oppdatert: 12. oktober 2026
\end{center}

\section{1. Hvem er ansvarlig?}
Skallstuggu Fjellstue AS er behandlingsansvarlig for opplysningene vi bruker til booking, kundekontakt og utleie. Org.nr. 919 472 669. Foretaksadresse: Tore Hunds gate 21, 7654 Verdal.

E-post: post@skallstuggu.no. Christian Green: 906 88 873. Monica Helden: 970 85 139.

Erklæringen gjelder bookingløsningen og nettsidene for Skallstuggu, inkludert skallstuggu-test.no.


\section{2. Opplysninger og kilder}
Vi behandler navn, e-postadresse, telefonnummer, ønsket eller avtalt leieperiode, valgte tillegg, pris, bookingstatus, betalingsstatus og meldinger knyttet til oppholdet. Ved fakturering behandler vi også nødvendige faktura- og kundeopplysninger.

Opplysningene kommer fra deg, fra hovedleietaker som bestiller for en gruppe, eller fra tidligere korrespondanse og vår bookingkalender i Google Kalender. Eksisterende reservasjoner kan overføres til den nye bookingløsningen.

Navn, e-post og telefon er obligatoriske i nettsidens bookingskjema, slik at vi kan behandle forespørselen og kontakte deg. Fritekst er valgfritt. Hvis du ikke ønsker å bruke skjemaet, kan du kontakte oss direkte.

Ikke skriv fødselsnummer, betalingskortopplysninger eller helseopplysninger i det vanlige skjemaet eller i bookingmeldinger. Behov for særskilt tilrettelegging avtales direkte med oss. Unngå navn på andre gjester når opplysningene ikke er nødvendige.


\section{3. Formål og behandlingsgrunnlag}
\begin{itemize}
\item Booking, kundedialog og gjennomføring av leie: nødvendige tiltak før avtaleinngåelse og oppfyllelse av avtale, jf. personvernforordningen artikkel 6 nr. 1 bokstav b. For kontaktpersoner som representerer en organisasjon eller en annen leietaker bruker vi bokstav f; vår berettigede interesse er å administrere leieforholdet.
\item Fakturering og pliktig regnskapsdokumentasjon: rettslig plikt etter bokføringsreglene, jf. artikkel 6 nr. 1 bokstav c.
\item Teknisk drift, tilgangskontroll, forebygging av misbruk og håndtering av rettskrav: berettiget interesse, jf. artikkel 6 nr. 1 bokstav f. Interessen er å beskytte tjenesten, opplysningene og partenes krav.
\end{itemize}
Vi bruker ikke bookingopplysninger til markedsføring uten et eget gyldig grunnlag. Et eventuelt samtykke kan trekkes tilbake uten at tidligere lovlig behandling blir ulovlig. Du kan be om mer informasjon om våre interesseavveininger.


\section{4. Kalender og privat bookinglenke}
Den offentlige kalenderen viser opptatte datoer, ikke kundenavn, telefonnummer, e-post eller meldinger. Bookingopplysninger er tilgjengelige for autorisert administrasjon.

Den private bookinglenken gir tilgang til opplysninger om bestillingen og tilhørende samtale. Ikke del lenken med uvedkommende. Kontakt oss dersom den kommer på avveie. Bookingforespørsler behandles av oss; vi bruker ikke profilering eller helautomatiske avgjørelser med rettsvirkning for deg.


\section{5. Tjenester og mottakere}
\begin{itemize}
\item Supabase: lagring og behandling av bookingopplysninger, meldinger, tilgang og tekniske funksjoner. Hoveddatabasen er plassert i regionen eu-west-1 i Irland.
\item Resend: utsending av bookingrelaterte e-poster. Tjenesten mottar nødvendig mottakeradresse og innhold i e-posten, inkludert eventuell privat bookinglenke.
\item Fiken: kunde- og fakturahåndtering. Nødvendige kunde- og fakturaopplysninger registreres der når vi bruker tjenesten for bestillingen. Testløsningen er koblet til et testforetak.
\item Google Kalender og våre e-posttjenester: administrasjon av reservasjoner og korrespondanse. Relevante navn, datoer og kundenotater kan inngå.
\item GitHub Pages og leverandører som leverer nettsidens tekniske filer: publisering og levering av nettsiden. Ved besøk mottar disse tekniske opplysninger som IP-adresse og forespurt ressurs.
\end{itemize}
Opplysninger gis bare til medarbeidere og leverandører med tjenstlig behov, eller når vi er pålagt å utlevere dem. Vi selger ikke personopplysninger. Databehandlere skal behandle opplysninger etter våre instrukser og gjeldende databehandlervilkår.


\section{6. Behandling utenfor EØS}
Selv om hoveddatabasen ligger i Irland, bruker vi internasjonale leverandører. E-postlevering, drift, støtte og underleverandører kan innebære behandling utenfor EØS, blant annet i USA. Supabase har også virksomhet i Singapore.

Leverandørenes databehandlervilkår beskriver internasjonale overføringer og bruk av EUs standard personvernbestemmelser. Overføring skal ha et gyldig grunnlag etter personvernforordningen kapittel V, for eksempel standardbestemmelser med nødvendige tilleggstiltak eller en relevant beslutning om tilstrekkelig beskyttelsesnivå. Kontakt oss for informasjon om hvilke garantier som gjelder for dine opplysninger og for kopi eller lenke til dem.

Leverandørinformasjon: https://supabase.com/legal/dpa og https://resend.com/legal/dpa.


\section{7. Lagring og sletting}
Lagringstiden avhenger av om henvendelsen blir en avtale, om opplysninger inngår i regnskapet, og om det finnes et uavklart krav. Følgende frister gjelder som vår sletterutine:

\begin{itemize}
\item Forespørsler som ikke fører til bestilling: slettes senest 12 måneder etter at saken er avsluttet.
\item Bookingopplysninger og nødvendig avtaledialog: oppbevares normalt i inntil tre år etter avsluttet opphold og oppgjør, for oppfølging av avtalen og eventuelle krav. Fritekst som ikke lenger trengs slettes tidligere.
\item Fakturaer og annet pliktig regnskapsmateriale: oppbevares som hovedregel i fem år etter regnskapsårets slutt. Enkelte dokumenttyper kan ha en annen lovbestemt frist.
\item Ved en konkret tvist eller rettslig plikt beholdes relevante opplysninger til saken er avsluttet og nødvendig frist er utløpt. Øvrige opplysninger slettes etter normal rutine.
\end{itemize}
Sletting gjennomføres ved gjennomgang av bookingløsningen, kalenderen, e-post og kunderegisteret. Fristene betyr ikke at nettsiden allerede har automatisk sletting. Sikkerhetskopier og tekniske logger følger leverandørenes dokumenterte lagringstider og skal ikke brukes til andre formål. Du kan be om nærmere informasjon om disse tidene.


\section{8. Nettbesøk, kart og teknisk lagring}
Nettsiden bruker tekniske tjenester for å levere innhold og behandle forespørsler. Disse mottar blant annet IP-adresse, tidspunkt og tekniske forespørsler. Bookingløsningen begrenser antall forespørsler for å motvirke misbruk. Administratorinnlogging bruker lokal nettleserlagring for å holde administratoren innlogget.

På kontaktsiden ligger et innebygd Google-kart. Når kartet lastes, sendes tekniske opplysninger til Google, som behandler dem etter sine egne personvernvilkår. Se https://policies.google.com/privacy. Vi bruker ikke kartet til å motta din posisjon.

Erklæringen gir ikke samtykke til analyse- eller markedsføringssporing. Dersom slike funksjoner tas i bruk, må informasjon og eventuelt samtykke håndteres før de aktiveres.


\section{9. Dine rettigheter}
Du kan be om innsyn, kopi og retting av opplysninger om deg. Du kan også be om sletting eller begrensning når vilkårene er oppfylt. Opplysninger som vi plikter å oppbevare, for eksempel regnskapsbilag, kan ikke alltid slettes straks.

For opplysninger du har gitt oss og som behandles automatisk på grunnlag av avtale eller samtykke, kan retten til dataportabilitet gjelde.

Du har rett til å protestere mot behandling basert på berettiget interesse. Vi må da vurdere protesten og kan bare fortsette når regelverket gir adgang til det.

Send forespørselen til post@skallstuggu.no. Vi svarer normalt innen én måned. Vi kan be om nødvendig bekreftelse på identitet for å unngå utlevering til feil person. Du kan klage til Datatilsynet, se https://www.datatilsynet.no.


\section{10. Sikkerhet og endringer}
Vi bruker tilgangskontroll og kryptert forbindelse til bookingløsningen. Bare autoriserte brukere skal ha tilgang til administrasjonen. Del ikke private bookinglenker eller innloggingsopplysninger.

Ved vesentlige endringer i hvordan opplysningene brukes, gir vi berørte personer informasjon på egnet måte. Nye formål vurderes før behandlingen starter. Datoen på første side viser siste oppdatering av teksten.

\end{document}

